% Volume X: Normative Conclusions
% Part XVIII. Freedom Under Causal Opacity
% Chapter 107: Equality of Evidentiary Standing
% Spine: proof-spines/ch107-spine.md

\chapter{Equality of Evidentiary Standing}
\label{ch:ch107}

Formal equality is inadequate when different actors possess radically unequal ability to collect, preserve, interpret, and present evidence. Equality of evidentiary standing concerns practical access to the operations by which institutional reality is produced: who can get a record made, who can keep it alive, who can have it believed, and who can force it into the path of decision.

The issue is therefore not simply whether everyone may speak, but whether everyone can participate in the evidentiary machinery that turns speech into institutional fact. A free order requires more than neutral rules on paper; it requires conditions under which persons and groups are not structurally disabled from making their evidence count.
